\section{Industrial Case Study. Structural Biology Pipelines (RQ5)}
\label{sec:practice}
Silent-fix research motivates retaining repair purpose beyond code~\cite{colefunda}. We evaluate maintenance independently of SusVibes through ten feature tasks in five scientific repositories, separating repair memory from boundary feedback.

\begin{table*}[t]
\caption{Independent biology maintenance study: two tasks per repository. Repository rows report final FuncPass/SecPass/JointPass after at most one retry; aggregate first and final outcomes follow. FuncPass requires ordinary and feature checks; SecPass requires saved and extended boundary checks, including reliability contracts for DockQ and AutoDock-Vina. Totals have denominator ten.}
\label{tab:bio}
\centering\footnotesize
% Evidence: silent_fix_feature_v1/{cases,continuations}; repository counts:
% paper/NOSTRUM_BARCELONA_PROOFREAD_PACK.md, original pinned scan inventory.
\setlength{\tabcolsep}{3pt}
\renewcommand{\arraystretch}{1.12}
\begin{tabularx}{\textwidth}{@{}l r >{\raggedright\arraybackslash}p{0.17\textwidth} >{\raggedright\arraybackslash}X >{\raggedright\arraybackslash}X c c@{}}
\toprule
Repository & Alerts & Target symbol & Repaired component defect & Feature evolution & \multicolumn{2}{c}{Joint success}\\
\cmidrule(lr){6-7}
 & & & & & No Memory & Memory+Gate\\
\midrule
DockQ & 10 & \code{load_PDB}
& Handle leak on parser failure/fallback
& Extension-first CIF/PDB parsing, optional format return
& Yes & Yes\\
haddock3 & 247 & \code{extract_files_flat}
& Extraction escape through destination symlink
& ZIP and TAR support, optional nested paths
& No$^{a}$ & Yes\\
ProteinMPNN & 43 & \code{StructureDataset}\newline\code{.__init__}
& Unvalidated record names used in output paths
& Minimum-length and missing-coordinate filters
& Yes & Yes\\
RFdiffusion & 49 & \code{Sampler.}\newline\code{load_checkpoint}
& Unrestricted pickle loading under pinned runtime$^{b}$
& Legacy-checkpoint exception handling and CPU fallback
& No$^{c}$ & Yes, retry\\
AutoDock-Vina & 60 & \code{Vina.write_poses}
& Overwrite branch bypasses numeric validation
& Verbose flag and consolidated validation chain
& Yes & Yes\\
\midrule
\multicolumn{5}{@{}l}{\textbf{Joint successes, first proposal}} & \textbf{3/5} & \textbf{4/5}\\
\multicolumn{5}{@{}l}{\textbf{Joint successes, final accepted feature implementation}} & \textbf{3/5} & \textbf{5/5}\\
\multicolumn{5}{@{}l}{Additional proposals / secure-code fallbacks} & 0 / 0 & 1 / 0\\
\bottomrule
\end{tabularx}
\par\smallskip
\begin{minipage}{\textwidth}\scriptsize
$^{a}$ ZIP output is missing. The saved boundary check passes.
$^{b}$ User-specified remedy, locally implemented after two unsuccessful model repairs.
$^{c}$ An unimported \code{pickle} name leaves the feature and boundary outcomes unknown. The gate's retry resolves this in Memory+Gate.
\end{minipage}

\end{table*}

\subsection{Repair Contracts and Independent Scope}
With input from our Nostrum Biodiscovery collaborator, we selected RFdiffusion~\cite{rfdiffusion}, ProteinMPNN~\cite{proteinmpnn}, haddock3~\cite{haddock3}, DockQ~\cite{dockq}, AutoDock-Vina~\cite{vina}, and foldseek~\cite{foldseek}. Semgrep 1.101.0 reports 420 warnings. Source and history review identify five Python repair targets outside those warning locations. Foldseek contributes eleven warnings but no suitable project-owned Python case.

RFdiffusion restricts checkpoint loading with \code{weights_only=True} and \code{safe_globals([slice])} under PyTorch 2.5.1~\cite{pytorchserialization,pythonpickle}. Haddock3 checks resolved extraction destinations; ProteinMPNN restricts record names used in output paths. DockQ closes handles on parser failure, and AutoDock-Vina validates numeric arguments across overwrite branches. These last two are reliability contracts. Four repairs were model-written; RFdiffusion used a developer-specified remedy after two unsuccessful model attempts.

Each original boundary assertion fails before repair and passes afterward with an ordinary-input control. External JSON notes record invariants, rationale, source identity, and checks. We strip comments and docstrings with \code{tokenize} and \code{ast}, preserve the executable AST excluding documentation nodes, and reverify the cleaned starts~\cite{pythontokenize,pythonast}.

\subsection{Equal-Budget Feature Evolution}
The tasks extend behavior near these contracts (Table~\ref{tab:bio}), including legacy checkpoint compatibility, archive aliases, hierarchical names, in-memory pose output, and lazy parsing. This purposive stress sample has ten tasks but five repository clusters. No SusVibes continuation enters the evaluation.

We retain each default-mode check and add a check for the new mode: an approved hierarchy or Namespace must not fail merely because the old contract excluded it. Before generation, each repaired start passes its old checks; an evaluator-written reference passes all checks; and a weakened reference retains the feature but triggers a boundary assertion. Reference code is withheld from the model.

An oracle audit removed undocumented requirements to retain Namespace representation and avoid consuming model numbers before parsing. All four conditions of the three affected tasks were rerun. An adapter audit allowed the supplied \code{FilePath = Path} alias without changing generated code and replayed all four archive-alias conditions. The comparison retains 24 unchanged cells, twelve reruns, and four replays. Original runs are archived; prompts, security assertions, and retry budgets were unchanged.

DeepSeek-V3.2 at temperature zero receives the cleaned source and feature request in all four conditions:
\begin{itemize}
\item \textbf{B0: Functional gate.} Application, ordinary, and feature failures can trigger one retry.
\item \textbf{B1: Memory.} B0 plus the repair note and a prospective clarification of the extended contract.
\item \textbf{B2: Boundary gate.} B0 plus saved and extended boundary checks in retry feedback.
\item \textbf{B3: Both.} Memory and boundary feedback.
\end{itemize}
All cells have the same token budget and maximum of two proposals. B0/B1 boundary checks execute only after generation ends; their outcomes cannot affect retries. Checks use pinned offline containers with read-only source and evaluator mounts. All 40 cells are retained, without outcome-dependent stopping. The earlier unequal-budget comparison remains separate archival evidence.

FuncPass requires ordinary and feature checks; SecPass requires saved and extended boundary checks; JointPass requires both. We report first and final proposals separately. Execution errors are unknown boundary outcomes, not violations. A rejected final proposal leaves the repaired start in place; this fallback is not feature delivery.

Table~\ref{tab:bio} separates first and final outcomes. B0 accepts three violating proposals: one unrestricted checkpoint load and two portable backslash-name violations. The latter are contract failures, not demonstrated Linux path escapes. B1 accepts one proposal that relaxes the default simple-name restriction. B2 and B3 accept no observed violation, with seven and one fallbacks, respectively. B0 has one unknown outcome.

The five additional joint successes of B3 over B0 comprise one restricted-loading success, one portable-naming success, and three feature or execution fixes. Figure~\ref{fig:bio-memory} shows the deserialization contrast. The combined condition has the highest joint count in this cohort; five selected repository clusters do not establish a population effect.

\begin{figure}[t]
\centering
% Code excerpt from silent_fix_feature_v1/cases/haddock_extract_symlink/starting_code.py.
% The repair-note fields below are paraphrased for display, not literal JSON.
\fbox{\begin{minipage}{0.94\columnwidth}
\footnotesize
\textbf{1. Silent Security Repair in Code}\par
Comments stripped.\par\smallskip
{\ttfamily
 target\_path = dest / member.name\par
 target\_resolved = target\_path.resolve()\par
 if dest not in target\_resolved.parents \textbackslash\par
 \hspace*{1em}and dest != target\_resolved:\par
 \hspace*{2em}continue\par}
\smallskip\hrule\smallskip
\textbf{2. TRACER External Repair Note}\par
Summary of the saved JSON note.\par
\textbf{Target symbol.} \code{extract_files_flat}\par
\textbf{Security rule.} Final extraction paths must stay inside the destination directory after resolving symlinks.\par
\textbf{Linked boundary check.} Reject a symlink escape and accept safe archive members.\par
\smallskip\hrule\smallskip
\textbf{3. Subsequent Feature Evolution}\par
\textbf{Request.} Add ZIP/TAR and nested-path support.\par
\textbf{No Memory.} Fails the feature check.\par
\textbf{Memory + Gate.} Passes ordinary behavior, the new feature check, and the saved security boundary check.
\end{minipage}}

\caption{An observed RFdiffusion regression under equal retry budgets. Both final proposals pass ordinary and legacy-feature checks; B0 enables unrestricted deserialization, while B3 retains restricted loading and passes both boundary checks.}
\label{fig:bio-memory}
\end{figure}

\rqfinding{RQ5}{Can repair memory and boundary gates preserve repaired properties while subsequent feature changes succeed?}{With the same one-retry limit, joint success is 4/10 with functional feedback alone, 8/10 with memory, 3/10 with a boundary gate, and 9/10 with both. B3 rejects its remaining contract-violating proposal; fallback is not feature success.}
